\documentclass{SibFU-docs}

% доп. пакеты
\usepackage{graphicx}
\usepackage{float}
\usepackage{hyperref}
\usepackage{caption}
\usepackage{subcaption}
\usepackage{amsmath}
% предотвартим ошибку с \Bbbk перед загрузкой пакета amssymb
\let\Bbbk\relax
\usepackage{amssymb}
\usepackage{booktabs}
\usepackage{tabularx}
\usepackage{longtable}
\usepackage{enumitem}
\usepackage{multirow}

% гиперссылки
\hypersetup{
    colorlinks=true,
    linkcolor=black,
    citecolor=blue,
    urlcolor=blue
}


% библиография
\addbibresource{report.bib}

% титульный лист
\begin{document}
\setlist[itemize]{left=1.3cm, itemsep=1pt, topsep=2pt} % для выравнивания списков
\setlist[enumerate]{left=1.3cm, itemsep=1pt, topsep=2pt}
\renewcommand{\contentsname}{}
\mycovertitle
{Гуманитарный институт}
{Кафедра прикладной информатики в искусстве и интерактивном медиа}
{Технологии создания и обработки текстовых документов}
{ст. преподаватель И.~Р.~Нигматуллин}
{Захаров И.~М., Тетерина П.~А., Логинова Ю.~В}
 
\begin{center}
\bfseries РЕФЕРАТ
\end{center}
\noindent

Реферат по теме: «Технологии создания и обработки текстовых документов» содержит 47 страниц и 85 использованных источников.

Ключевые слова: КОДИРОВАНИЕ ТЕКСТА, UNICODE, UTF-8, ОТКРЫТЫЕ И ПРОПРИЕТАРНЫЕ ФОРМАТЫ, ЯЗЫКИ РАЗМЕТКИ (HTML, MARKDOWN, XML, LaTeX), СЕМАНТИЧЕСКАЯ ОРГАНИЗАЦИЯ, МЕТАДАННЫЕ (DUBLIN CORE), ОБРАБОТКА ЕСТЕСТВЕННОГО ЯЗЫКА (NLP), АВТОМАТИЧЕСКОЕ РЕФЕРИРОВАНИЕ, АНАЛИЗ ТОНАЛЬНОСТИ, ДОЛГОВРЕМЕННОЕ ХРАНЕНИЕ, АРХИВАЦИЯ, PDF/A, ЦЕЛОСТНОСТЬ ДАННЫХ, СТРАТЕГИИ СОХРАНЕНИЯ (МИГРАЦИЯ, ЭМУЛЯЦИЯ)

В работе проводится комплексное исследование современных технологий и методологий работы с текстовыми документами — от их фундаментального представления в цифровой форме до продвинутого анализа, семантической организации и долговременного архивного хранения.

Рассмотрены эволюция и принципы текстового кодирования — от базового ASCII к универсальному стандарту Unicode и эффективным кодировкам семейства UTF, обеспечивающим корректное представление многоязычных текстов. Детально проанализированы особенности проприетарных и открытых форматов текстовых файлов, их преимущества, недостатки и области применения, а также роль языков разметки (HTML, Markdown, XML, LaTeX) в создании структурированных, машиночитаемых документов. Особое внимание уделено принципам семантической организации текста и использованию метаданных (включая стандарт Dublin Core) для улучшения поиска, обработки и долговременной сохранности документов.

Проанализированы современные технологии автоматической обработки и анализа текстовой информации: от предобработки (токенизация, лемматизация) до распознавания именованных сущностей, синтаксического и семантического анализа, анализа тональности, тематического моделирования и автоматического реферирования. Завершающая часть работы посвящена основам долговременного хранения и архивации текстовых ресурсов, включая стратегии сохранения (миграция, эмуляция), форматы долговременного хранения (PDF/A, XML), методы обеспечения целостности данных и баланс между сохраняемостью и доступностью цифровых объектов.

Основные выводы:
\begin{enumerate}
\item Универсальный стандарт Unicode и эффективные кодировки (UTF-8) являются фундаментом корректного представления, обработки и обмена многоязычными текстовыми данными в цифровой среде, обеспечивая совместимость и устойчивость к ошибкам;
\item Выбор формата текстового файла должен быть осознанным: открытые стандарты (ODT, PDF/A, XML) гарантируют долговременную доступность и независимость от вендоров, тогда как проприетарные форматы могут предоставлять расширенные возможности в рамках конкретных экосистем;
\item Языки разметки и семантическая организация текста отделяют содержание от оформления, превращая документы в структурированные, машиночитаемые объекты, что значительно расширяет возможности их автоматической обработки, преобразования и интеграции;
\item Современные технологии обработки естественного языка (NLP) на основе нейросетевых моделей позволяют автоматически извлекать знания из текста, анализировать тональность, определять темы и генерировать краткие содержания, открывая новые возможности для работы с большими текстовыми массивами;
\item Долговременное сохранение цифровых текстовых ресурсов требует комплексного подхода, сочетающего технические стратегии (миграция, эмуляция), использование специализированных форматов хранения, обеспечение целостности данных и ведение подробных метаданных, что обеспечивает доступность цифрового наследия для будущих поколений.
\end{enumerate}

\vspace{1em}

Рекомендации:
\begin{itemize}[label=-]
\item При создании текстовых документов всегда использовать кодировку UTF-8 для обеспечения корректного отображения символов всех языков и избежания проблем совместимости;
\item Для долговременного хранения и архивации важных документов отдавать предпочтение открытым, стандартизированным форматам, таким как PDF/A (для фиксированного представления) и XML/TXT (для структурированных данных), и избегать хранения критичной информации исключительно в проприетарных форматах;
\item Активно использовать языки разметки (Markdown для простых документов, LaTeX для научных работ, HTML для веб-контента) и семантические принципы организации текста, что облегчает автоматическую обработку, преобразование и долговременную сохранность документов;
\item Внедрять метаданные (например, по стандарту Dublin Core) для описания текстовых ресурсов, что значительно улучшает их поиск, каталогизацию и управление как в личных, так и в корпоративных и архивных системах;
\item Осваивать базовые методы и инструменты автоматической обработки текста (токенизация, стемминг, анализ тональности) для эффективной работы с большими текстовыми массивами в исследовательской, аналитической и управленческой деятельности.
\end{itemize}

\newpage

\thispagestyle{empty}
\begin{center}
\bfseries СОДЕРЖАНИЕ
\end{center}
\vspace{-6pt}
\tableofcontents
\clearpage

\section*{}
\vspace*{-2.5em}
\begin{center}\textbf{ВВЕДЕНИЕ}\end{center}
\phantomsection
\addcontentsline{toc}{section}{Введение}

Текстовые документы являются одной из основных форм представления информации в цифровую эпоху. От простых заметок до сложных научных статей, от юридических документов до художественных произведений — тексты составляют основу коммуникации, образования, науки и бизнеса. Однако кажущаяся простота текстового документа скрывает за собой сложный комплекс технологий, обеспечивающих его создание, обработку, хранение и анализ.

Современные технологии работы с текстовыми документами охватывают широкий спектр вопросов: от низкоуровневого представления символов в цифровой форме до высокоуровневых методов автоматического понимания смысла текста. Развитие стандартов кодирования, появление открытых форматов файлов, распространение языков разметки, внедрение семантических принципов организации информации, прогресс в области обработки естественного языка и разработка стратегий долговременного хранения — все эти направления формируют целостную экосистему работы с текстовой информацией.

Актуальность темы обусловлена экспоненциальным ростом объёмов текстовой информации и повышением требований к её качеству, доступности и долговременной сохранности. Грамотное владение технологиями создания и обработки текстовых документов становится критически важной компетенцией для специалистов в различных областях.

Цель работы: исследовать совокупность технологий и методов, составляющих современный процесс создания, обработки, анализа и хранения текстовых документов.

Задачи:
\begin{enumerate}
\item Изучить принципы представления и кодирования текстовых данных в цифровой форме;
\item Проанализировать особенности проприетарных и открытых форматов текстовых файлов;
\item Рассмотреть роль языков разметки в структурировании документов;
\item Исследовать принципы семантической организации и использования метаданных в текстовых файлах;
\item Проанализировать современные технологии автоматической обработки и анализа текстовой информации;
\item Изучить основы долговременного хранения и архивации текстовых ресурсов.
\end{enumerate}

Методы: системный анализ, сравнительный анализ технологий и стандартов, обобщение и структурирование информации из профессиональных источников.

\newpage
% вручную увеличить номер секции и записать в оглавление с номером,
% чтобы в теле не печатался автоматический номер, но в tableofcontents он остался
\refstepcounter{section}
\addcontentsline{toc}{section}{\protect\numberline{\thesection} Представление и кодирование текстовых данных в цифровой форме}
\vspace*{-2.5em}
\begin{center}\textbf{\thesection\ ПРЕДСТАВЛЕНИЕ И КОДИРОВАНИЕ ТЕКСТОВЫХ ДАННЫХ В ЦИФРОВОЙ ФОРМЕ}\end{center}
\vspace*{-1.5em}

\ManualSubsection{Фундаментальные принципы цифрового представления текста}

В основе всех современных вычислительных систем лежит принцип представления информации в виде двоичных данных — последовательностей нулей и единиц. Текстовая информация, как и любая другая, должна быть преобразована в эту форму для хранения и обработки компьютером. Процесс преобразования символов — букв, цифр, знаков препинания — в числовые коды, а затем в двоичные последовательности, называется кодированием. Обратный процесс восстановления символов из двоичного представления именуется декодированием \cite{unicode2024, petzold2001}.

Каждый текст состоит из дискретных единиц — символов. Задача кодирования заключается в сопоставлении каждому символу уникального числового идентификатора, называемого кодовой точкой. Совокупность всех возможных кодовых точек образует кодовое пространство. Алгоритм, определяющий, как набор кодовых точек преобразуется в последовательность байтов для хранения или передачи, называется кодировкой или кодовой страницей \cite{knut2020, tanenbaum2015}.

Базовый процесс кодирования включает несколько этапов: выбор символа для сохранения, сопоставление его с кодовой точкой согласно таблице кодирования, преобразование числового значения в двоичную последовательность и сохранение полученных битов в файле или передача по сети. Декодирование представляет собой обратную последовательность действий: чтение битов, определение кодовой точки и отображение соответствующего символа \cite{rfc3629, w3c2024}.

\ManualSubsection{Историческая эволюция стандартов кодирования}

История текстового кодирования отражает развитие вычислительной техники и рост потребностей в обработке многоязычной информации. Первым широко распространённым стандартом стал ASCII (American Standard Code for Information Interchange), разработанный в 1960-х годах \cite{unicode1996}.

ASCII использует 7 бит для кодирования, что позволяет представить 128 уникальных символов. Диапазон от 0 до 31 отведён под управляющие символы, такие как перевод строки (LF) и возврат каретки (CR). Диапазон от 32 до 126 содержит печатные символы: пробел, цифры, заглавные и строчные латинские буквы, основные знаки препинания. Главным ограничением ASCII стала ориентация исключительно на американский английский, что сделало невозможным представление букв других языков \cite{tannenbaum2019}.

Для преодоления этого ограничения стали использоваться 8-битные кодировки, которые расширили кодовое пространство до 256 символов. Первые 128 позиций оставались идентичными ASCII, а диапазон 128–255 использовался для национальных символов. Появилось семейство стандартов ISO/IEC 8859, где каждая часть отвечала за определённую группу языков: ISO-8859-1 для западноевропейских языков, ISO-8859-5 для кириллицы. Параллельно развивались национальные стандарты, такие как KOI8-R для русского языка и Windows-1251 от Microsoft \cite{iana2024}.

Основной проблемой 8-битных кодировок стала их несовместимость: текст, созданный в одной кодировке, при открытии в другой отображался как бессмысленный набор символов. Эта проблема, известная как «кракозябры», серьёзно ограничивала международный обмен данными и требовала универсального решения \cite{iso10646}.

\ManualSubsection{Современный стандарт Unicode и семейство UTF}

Универсальным решением проблемы множества кодировок стал стандарт Unicode, основная идея которого заключается в создании единой таблицы для всех символов всех письменностей мира \cite{unicode2024}.

Unicode основывается на нескольких ключевых принципах. Принцип универсальности предполагает один стандарт для всех языков и платформ. Принцип однозначности гарантирует, что каждому символу присваивается уникальная кодовая точка. Принцип ёмкости обеспечивает поддержку более 1,1 миллиона кодовых точек, что достаточно для всех известных и будущих письменностей. Кодовая точка в Unicode записывается в формате U+XXXX, где XXXX — шестнадцатеричное число \cite{unicode1996}.

Важно понимать, что Unicode определяет только соответствие «символ → кодовая точка», но не задаёт, как эти кодовые точки будут представлены в виде байтов. Для этого используются алгоритмы кодирования, известные как Unicode Transformation Format (UTF) \cite{rfc3629}.

UTF-32 кодирует каждую кодовую точку Unicode ровно одним 32-битным числом. Это обеспечивает простоту и скорость доступа к произвольному символу, но приводит к чрезмерному расходу памяти, что ограничивает практическое применение этого формата \cite{w3c2024}.

UTF-16 использует 16-битные слова для представления символов. Базовые символы кодируются одним словом, для остальных применяется система суррогатных пар. Этот формат широко используется внутри операционных систем Windows и в таких технологиях, как Java и .NET, но имеет недостаток в виде непостоянной длины символа \cite{tanenbaum2015}.

UTF-8 стал стандартом де-факто для веба, файловых систем и систем хранения данных. Он использует схему с переменной длиной от 1 до 4 байт. Символы ASCII кодируются одним байтом и полностью совпадают с оригинальным стандартом, что обеспечивает обратную совместимость. Для остальных символов используются 2, 3 или 4 байта \cite{rfc3629}.

Преимущества UTF-8 включают обратную совместимость с ASCII, эффективное использование памяти для западноевропейских языков, устойчивость к повреждениям (повреждение одного байта не разрушает всю последующую строку) и независимость от порядка байтов (endianness) \cite{w3c2024}.

\ManualSubsection{Практические аспекты и проблемы кодирования}

На практике работа с текстовым кодированием сопряжена с несколькими важными аспектами. Первый из них — определение кодировки файла. Компьютер не «видит» кодировку автоматически; для корректного отображения текста он должен знать, какой алгоритм декодирования применять \cite{unicode2024}.

Для указания кодировки могут использоваться различные механизмы. В вебе стандартом является UTF-8, и серверы обычно передают информацию о кодировке в HTTP-заголовках. В файловых системах иногда применяется сигнатура — специальный непечатаемый символ U+FEFF (Byte Order Mark, BOM), записываемый в начало файла. Однако использование BOM может вызывать проблемы совместимости, особенно в скриптах и исходных кодах \cite{iana2024}.

Другой важный аспект — нормализация текста. В Unicode одни и те же визуальные символы могут быть представлены разными последовательностями кодовых точек. Например, буква «ё» может быть записана как одна кодовая точка U+0451 или как композиционная последовательность: U+0435 («е») + U+0308 (комбинирующий диерезис). Нормализация — это процесс приведения текста к стандартной, канонической форме для обеспечения корректного сравнения, поиска и обработки \cite{w3c2024}.

Современные операционные системы и приложения должны корректно обрабатывать Unicode во всех его аспектах. Это включает поддержку различных направлений письма (слева направо, справа налево, сверху вниз), правильное отображение комбинирующих символов, обработку суррогатных пар и многое другое \cite{tanenbaum2015}.

Понимание принципов кодирования текста критически важно для разработчиков, системных администраторов, специалистов по обработке данных и всех, кто работает с многоязычной цифровой информацией. Это знание позволяет избежать распространённых ошибок, обеспечивает корректную обработку текстовых данных и способствует созданию интернационализированных приложений и систем \cite{knut2020, petzold2001}.

\newpage
\refstepcounter{section}
\addcontentsline{toc}{section}{\protect\numberline{\thesection}Особенности проприетарных и открытых форматов текстовых файлов}
\begin{center}\textbf{\thesection\ ОСОБЕННОСТИ ПРОПРИЕТАРНЫХ И ОТКРЫТЫХ ФОРМАТОВ ТЕКСТОВЫХ ФАЙЛОВ}\end{center}
\vspace*{-1.5em}

\ManualSubsection{Классификация текстовых файлов и основные понятия}

В цифровую эпоху текст сохраняется и обрабатывается в различных форматах файлов, которые можно разделить на две фундаментальные категории: проприетарные (закрытые) и открытые. Понимание различий между ними критически важно для обеспечения долговременной сохранности данных, совместимости и цифровой автономии \cite{kyrjakov2018, semenov2019}.

Текстовый файл представляет собой файл, содержащий данные в виде последовательности символов, закодированных в соответствии с определённой кодировкой, такой как UTF-8 или Windows-1251. Эти файлы предназначены для чтения человеком и обработки текстовыми редакторами \cite{odf2024}.

Существуют два основных типа текстовых файлов. Простой текст (plain text) содержит только последовательность символов, информацию о переносах строк и базовое форматирование, такое как табуляция и пробелы. Он не включает информацию о шрифтах, цветах или размерах. Примерами являются файлы с расширениями .txt, .csv, а также исходные коды программ. Форматированный текст (rich text) содержит не только текст, но и метаданные о его внешнем виде: шрифты, стили, выравнивание, изображения, таблицы. Примерами служат форматы .docx, .odt, .pdf, .rtf \cite{iso26300, ecma376}.

\ManualSubsection{Проприетарные форматы: характеристики, преимущества и недостатки}

Проприетарный формат — это формат, спецификации которого являются собственностью компании-разработчика. Его внутренняя структура закрыта, не документирована для публики или защищена патентами, что ограничивает возможности независимой реализации поддержки такого формата \cite{msft2024}.

Классическими представителями проприетарных форматов являются бинарные форматы Microsoft Office до версии 2007: .doc для текстовых документов, .xls для электронных таблиц, .ppt для презентаций. Эти форматы отличаются сложной бинарной структурой (Binary Interchange File Format, BIFF), которая обеспечивала высокую производительность в родных приложениях, но создавала серьёзные проблемы с совместимостью и безопасностью \cite{petrov2022}.

Современные форматы Microsoft Office — .docx, .xlsx, .pptx — основаны на стандарте Office Open XML (OOXML), стандартизированном как ECMA-376 и ISO/IEC 29500. Несмотря на название «Open XML», эти форматы сохраняют элементы проприетарности, так как контроль над спецификацией остаётся за Microsoft, а полная реализация всех возможностей требует лицензирования \cite{ecma376, msft2024}.

Преимущества проприетарных форматов включают богатую функциональность, часто превосходящую возможности открытых аналогов. Они обеспечивают глубокую интеграцию с экосистемой продуктов компании-разработчика, что создаёт удобство для пользователей, работающих исключительно в этой среде. Для рядового пользователя в рамках замкнутой экосистемы такие форматы обеспечивают простоту и удобство работы \cite{gnu2024}.

Однако недостатки проприетарных форматов существенны. Главный из них — зависимость от вендора: для полного доступа к функционалу файла требуется конкретное программное обеспечение. Альтернативные программы могут некорректно отображать сложное форматирование. Проблемы совместимости проявляются, когда файлы, созданные в новых версиях программы, не открываются или неправильно отображаются в старых версиях. Закрытость спецификаций затрудняет разработку стороннего ПО, что приводит к ошибкам и ограничениям в совместимости. Сложная бинарная структура старых форматов и возможность запуска макросов делают их популярным вектором для кибератак. Наконец, существует риск для долговременного хранения: если компания-разработчик прекратит поддержку формата, доступ к данным может быть утрачен \cite{semenov2019, ivanova2021}.

\ManualSubsection{Открытые форматы: принципы, стандарты и практическое значение}

Открытый формат — это формат, спецификации которого публично доступны, бесплатны для использования и могут быть реализованы любым разработчиком в своём ПО без юридических или финансовых ограничений. Такие форматы обычно стандартизируются международными организациями, что гарантирует их стабильность и долговременную поддержку \cite{odf2024, iso26300}.

Яркими примерами открытых форматов являются простой текст (.txt), Comma-Separated Values (.csv) для табличных данных, OpenDocument Format (.odt, .ods, .odp), разработанный сообществом OASIS и принятый как международный стандарт ISO/IEC 26300. HyperText Markup Language (.html) как основа веба и Portable Document Format (.pdf), стандартизированный ISO, также относятся к открытым стандартам, хотя некоторые расширенные функции PDF могут оставаться проприетарными \cite{iso32000, w3c2024}.

Преимущества открытых форматов многообразны. Они обеспечивают независимость от вендора, так как отсутствует привязка к конкретному производителю ПО. Документ можно открыть в любой программе, поддерживающей данный формат. Благодаря открытым спецификациям разработчики со всего мира могут обеспечить корректную работу с форматом в своих приложениях на разных платформах. Долговременная сохранность гарантируется доступностью спецификаций: даже если все текущие программы перестанут существовать, в будущем можно будет написать новую программу для чтения этих файлов. Прозрачность и безопасность обеспечиваются простотой текстовых форматов, которые легко проверяемы на наличие вредоносного кода. Отсутствие скрытых макросов снижает поверхность атаки. Открытость стимулирует конкуренцию, позволяя создавать множество конкурирующих приложений, что ведёт к улучшению качества ПО и снижению цен \cite{gnu2024, voronov2020}.

Недостатки открытых форматов относительно невелики. Возможна потеря сложного форматирования при переходе из очень сложного проприетарного формата в открытый, так как некоторые продвинутые элементы могут быть упрощены или утеряны. Самые простые открытые форматы не поддерживают форматирования, что ограничивает их применение. Кроме того, хотя спецификация едина, разные программы могут по-своему интерпретировать её нюансы, что иногда приводит к мелким несовпадениям в отображении \cite{petrov2022}.

\ManualSubsection{Сравнительный анализ и рекомендации по выбору форматов}

Сравнительный анализ проприетарных и открытых форматов позволяет сформулировать рекомендации по их выбору для различных задач \cite{kyrjakov2018}.

Для хранения важных, долгосрочных данных, архивов, государственных документов предпочтительны открытые форматы, такие как ODT, PDF/A, TXT. Это гарантирует доступность информации в будущем независимо от изменений на рынке программного обеспечения.

Для совместной работы в гетерогенной среде, когда участники используют разные операционные системы и программное обеспечение, также лучше использовать открытые форматы (ODT) или стандартизированные (PDF). Это минимизирует проблемы с совместимостью и обеспечивает равный доступ всем участникам.

Для работы в рамках строго определённой корпоративной экосистемы, где все используют Microsoft 365 или аналогичные проприетарные решения, допустимо применение проприетарных форматов (.docx). Это обеспечит доступ ко всему функционалу и лучшую интеграцию.

Для обмена данными, которые должны быть легко обработаны программами (логи, датасеты, исходный код), оптимальны простые открытые форматы (.txt, .csv, .json, .xml). Их простота гарантирует лёгкость парсинга и обработки.

Для создания документов со сложным профессиональным форматированием (научные статьи, книги) в среде, где это форматирование должно быть точно воспроизведено, можно использовать PDF как открытый стандарт или DOCX с оговоркой на совместимость \cite{iso32000, ecma376}.

Проприетарные форматы предлагают мощный функционал и удобство в рамках замкнутой экосистемы, но за это приходится платить зависимостью от вендора и рисками для будущего. Открытые форматы являются залогом свободы, совместимости и долговечности данных, хотя иногда могут уступать в продвинутых функциях. Осознанный выбор формата, основанный на понимании этих компромиссов, является важным элементом цифровой грамотности \cite{gnu2024, voronov2020}.

\newpage
\refstepcounter{section}
\addcontentsline{toc}{section}{\protect\numberline{\thesection}Роль языков разметки в структурировании документа}
\begin{center}\textbf{\thesection\ РОЛЬ ЯЗЫКОВ РАЗМЕТКИ В СТРУКТУРИРОВАНИИ ДОКУМЕНТА}\end{center}
\vspace*{-1.5em}

\ManualSubsection{Основные понятия и назначение языков разметки}

Языки разметки представляют собой специальные способы оформления текста, которые помогают компьютеру понять структуру документа. Если обычный текст состоит просто из слов, то текст с разметкой содержит дополнительные указания о том, как эти слова должны располагаться и оформляться. Это позволяет программам корректно интерпретировать содержание документа и правильно его отображать или обрабатывать \cite{knuth1984texbook, grune2012markup}.

Самый известный пример языка разметки — HTML (HyperText Markup Language), на котором строятся веб-страницы. Когда пользователь видит на веб-странице заголовок, жирный текст или список, за всем этим стоит специальная разметка, которая указывает браузеру, как именно отображать эти элементы. Таким образом, языки разметки служат мостом между человеком, создающим документ, и компьютером, который этот документ обрабатывает \cite{bernerslee1992html, w3c_html}.

Основная функция языков разметки — отделение содержания документа от его оформления. Это означает, что автор указывает, какая часть текста является заголовком, какая — абзацем, где находится список или таблица, а программа уже решает, как эти элементы визуально представить. Такой подход делает работу с текстом более организованной и систематизированной, а также обеспечивает согласованность оформления во всём документе \cite{latex_project, xml_w3c}.

\ManualSubsection{Популярные языки разметки и их применение}

Современная практика работы с документами использует несколько широко распространённых языков разметки, каждый из которых имеет свои особенности и области применения \cite{markdown_spec, latex_project}.

HTML является основным языком для создания веб-страниц. С его помощью разработчики описывают структуру сайта: расположение заголовков, абзацев, изображений, ссылок, навигационных меню и других элементов. Браузер читает HTML-код и преобразует его в визуально привлекательную страницу. Современный HTML включает семантические элементы, такие как <header>, <article>, <section>, <footer>, которые не только определяют внешний вид, но и указывают смысловую роль каждого блока \cite{w3c_html}.

Markdown пользуется популярностью у программистов, технических писателей и пользователей, которым нужен простой способ создания форматированного текста. Его синтаксис интуитивно понятен: для заголовка достаточно поставить решётку, для жирного текста — двойные звёздочки, для списка — дефисы. Текст в Markdown легко читается даже в исходном виде, что делает его идеальным для документации, файлов README, сообщений на форумах и в системах обмена сообщениями \cite{markdown_spec}.

XML (eXtensible Markup Language) применяется, когда требуется организовать структурированный обмен данными между разными программами и системами. XML позволяет определять собственные теги и структуры, что делает его чрезвычайно гибким. Этот формат используется для выгрузки товаров из интернет-магазинов, передачи финансовой отчётности, описания конфигураций программного обеспечения и во многих других сценариях, где важна точность и строгость структуры \cite{xml_w3c}.

LaTeX незаменим для научных работ, где критически важны точность оформления математических формул, таблиц, графиков и библиографических ссылок. Многие научные журналы и издательства требуют подачи статей именно в формате LaTeX, так как он обеспечивает типографическое качество, недостижимое при использовании обычных текстовых процессоров \cite{knuth1984texbook, latex_project}.

\ManualSubsection{Преимущества использования языков разметки}

Использование языков разметки даёт несколько значительных преимуществ по сравнению с работой в простом тексте или проприетарных форматах документов \cite{grune2012markup, markdown_yaml}.

Структурированность документа означает, что каждый элемент имеет чёткое назначение. Заголовки отделены от основного текста, списки оформлены единообразно, таблицы имеют ясную структуру. Это особенно важно для больших документов, где необходимо быстро находить нужную информацию. Структурированные документы легче навигационно осваивать, проще автоматически обрабатывать и преобразовывать в другие форматы.

Переносимость обеспечивается тем, что документ, размеченный с помощью языка разметки, можно открывать на разных устройствах и в разных программах. Текст в Markdown будет одинаково хорошо выглядеть и в простом текстовом редакторе, и в специализированном приложении, и на веб-странице. Это устраняет проблему совместимости, характерную для проприетарных форматов.

Долговременная сохранность — ключевое преимущество текстовых форматов с разметкой. В то время как файлы, созданные в старых версиях проприетарных текстовых процессоров, могут перестать открываться в новых версиях, HTML и Markdown остаются читаемыми всегда. Простота и открытость этих форматов гарантируют, что даже через десятилетия можно будет написать программу для их обработки.

Автоматическая обработка становится возможной благодаря тому, что программы могут анализировать структуру документа, определённую разметкой. Можно автоматически генерировать оглавления, создавать индекс ключевых слов, преобразовывать документ в другие форматы (PDF, EPUB, HTML), извлекать определённые типы контента. Это значительно экономит время и снижает вероятность ошибок при ручной обработке.

\ManualSubsection{Практическое применение и влияние на совместную работу}

В современной цифровой среде языки разметки используются повсеместно, часто даже неявно для конечного пользователя \cite{bernerslee2001semanticweb, dublincore}.

Когда пользователь пишет сообщение в мессенджере и выделяет текст жирным или курсивом, он использует простейшую форму разметки. Когда автор форматирует текст в Google Docs или Microsoft Word, программа сохраняет это оформление с помощью внутренней разметки, основанной на XML. Для блогеров и контент-менеджеров знание HTML необходимо для тонкой настройки внешнего вида статей. Программисты используют Markdown для документации к своему коду. Учёные и студенты применяют LaTeX для написания курсовых, дипломных работ и научных статей.

Языки разметки играют важную роль в системах управления контентом и организации совместной работы над документами. В вики-системах, которые используются для создания баз знаний внутри компаний, сотрудники редактируют страницы с помощью упрощённой разметки. Это позволяет поддерживать единый стиль оформления даже когда над документами работают десятки людей.

В системах контроля версий, используемых программистами для совместной разработки, вся документация обычно ведётся в Markdown. Это обеспечивает консистентность и лёгкость сравнения разных версий документов. При создании технической документации часто используется подход «единый источник истины»: контент хранится в виде размеченного текста, а затем автоматически генерируется в разных форматах: PDF для печати, HTML для сайта, CHM для справки.

\ManualSubsection{Будущее языков разметки}

Развитие технологий продолжает влиять на эволюцию языков разметки. Появляются гибридные форматы, сочетающие простоту Markdown с мощностью XML. Развиваются инструменты визуального редактирования, которые скрывают сложность разметки от пользователя, но сохраняют её преимущества в виде чёткой структуры и machine-readability \cite{markdown_yaml, xml_w3c}.

Искусственный интеллект начинает использоваться для автоматического структурирования текста: программы учатся сами определять, где должны быть заголовки, списки, важные фрагменты. Однако базовые принципы языков разметки — разделение содержания и оформления, чёткая структура, machine-readability — остаются неизменными. Эти принципы обеспечивают долголетие и универсальность языков разметки в цифровую эпоху \cite{bernerslee2001semanticweb}.

Понимание и использование языков разметки становится важным навыком для эффективной работы с информацией в современном обществе. Освоение даже базовых принципов разметки позволяет создавать более качественные и организованные документы, упрощает совместную работу и обеспечивает надёжное хранение информации на долгие годы.

\newpage
\refstepcounter{section}
\addcontentsline{toc}{section}{\protect\numberline{\thesection}Принципы семантической организации и метаданных текстовых файлов}
\begin{center}\textbf{\thesection\ ПРИНЦИПЫ СЕМАНТИЧЕСКОЙ ОРГАНИЗАЦИИ И МЕТАДАННЫХ ТЕКСТОВЫХ ФАЙЛОВ}\end{center}
\vspace*{-1.5em}

\ManualSubsection{Семантическая организация текста: основные принципы и значение}

Семантическая организация текста представляет собой способ структурирования информации, при котором каждому элементу документа присваивается определённый смысл и назначение. В отличие от простого визуального форматирования, семантическая разметка указывает компьютеру не только как отображать текст, но и что означает каждая его часть — где находится заголовок, основной абзац, цитата, список или таблица \cite{bernerslee2001semanticweb, gilliland2016metadata}.

Такой подход позволяет программам лучше понимать содержание документа и корректно его обрабатывать. Например, поисковые системы могут точнее определять релевантность веб-страниц, анализируя семантические заголовки и разделы. Программы чтения для слабовидящих (скринридеры) правильно озвучивают структуру текста, сообщая пользователю, что сейчас читается заголовок, а что — обычный абзац. Системы автоматической обработки могут извлекать нужную информацию, опираясь на смысловые метки, а не на визуальные подсказки \cite{w3c_html, pdf_xmp}.

Семантическая разметка строится на нескольких ключевых принципах. Принцип смыслового разделения предполагает, что каждый элемент имеет чёткое назначение: заголовок, параграф, список, таблица. Принцип иерархической структуры устанавливает отношения вложенности между элементами: разделы содержат подразделы, списки содержат пункты. Принцип единообразия требует, чтобы одинаковые элементы оформлялись одинаково: все заголовки одного уровня выглядят единообразно. Принцип доступности гарантирует, что структура документа понятна людям с ограниченными возможностями и может быть корректно обработана вспомогательными технологиями \cite{dublincore, markdown_yaml}.

\ManualSubsection{Метаданные: понятие, типы и роль в работе с документами}

Метаданные — это информация о информации. В контексте текстовых файлов метаданные представляют собой дополнительные сведения о самом документе, которые не отображаются в основном содержимом, но критически важны для его обработки, поиска и управления \cite{gilliland2016metadata, dublincore}.

Метаданные могут включать широкий спектр информации: автора документа, дату создания, ключевые слова, язык текста, версию формата, права доступа, историю изменений и многое другое. Они помогают системам каталогизировать, искать и управлять документами без необходимости глубокого анализа всего их содержимого \cite{pdf_xmp, markdown_yaml}.

Современные текстовые форматы поддерживают различные типы метаданных, служащие разным целям. Описательные метаданные содержат информацию о содержании документа: автор, заголовок, ключевые слова, аннотация. Административные метаданные включают служебную информацию: дату создания, формат, размер файла, права доступа, историю версий. Структурные метаданные описывают строение документа: оглавление, связи между разделами, порядок страниц. Технические метаданные указывают параметры обработки: кодировку текста, версию формата, требования к программному обеспечению \cite{dublincore, gilliland2016metadata}.

Наличие полных и корректных метаданных значительно повышает ценность текстового документа, превращая его из простого файла в структурированный информационный объект, пригодный для автоматической обработки, долговременного хранения и эффективного поиска \cite{bernerslee2001semanticweb}.

\ManualSubsection{Практическая реализация в популярных форматах}

Разные форматы текстовых файлов по-разному реализуют принципы семантической организации и хранения метаданных, отражая их историю, назначение и философию разработки \cite{w3c_html, pdf_xmp}.

HTML использует систему тегов для семантической разметки. Современный HTML включает элементы вроде <header>, <article>, <section>, <footer>, которые чётко определяют назначение каждого блока. Метаданные размещаются в секции <head> документа и могут включать заголовок страницы, описание, ключевые слова, кодировку, информацию об авторе и многое другое. Стандарт Open Graph Protocol расширяет метаданные HTML для улучшения отображения страниц в социальных сетях \cite{w3c_html}.

Markdown, хотя и является упрощенным языком разметки, тоже поддерживает базовую семантику через символы выделения заголовков, списков, цитат. Метаданные могут добавляться через специальные блоки в начале документа в формате YAML (YAML Ain't Markup Language) или JSON. Такие блоки, часто называемые «front matter», содержат информацию о заголовке, авторе, дате, тегах и других атрибутах документа \cite{markdown_yaml}.

DOCX (формат Microsoft Word) хранит сложную систему метаданных, включающую информацию об авторе, организации, времени редактирования, статистику документа (количество слов, страниц, символов). Семантическая разметка реализована через стили, которые присваивают элементам определённую смысловую роль: «Заголовок 1», «Обычный текст», «Блок цитаты». Это позволяет отделить оформление от структуры и обеспечивает согласованность документа \cite{ecma376}.

PDF содержит обширные метаданные в блоке XMP (Extensible Metadata Platform), который может включать информацию о создателе, источнике, лицензии, ключевых словах. Современные PDF также поддерживают семантическую разметку через тегирование контента, что улучшает доступность для людей с ограниченными возможностями и точность поиска по документу. Стандарт PDF/A для долговременного архивирования требует наличия определённых метаданных для обеспечения сохранности и идентифицируемости документа \cite{pdf_xmp, iso32000}.

\ManualSubsection{Преимущества семантического подхода и использования метаданных}

Использование семантической организации и метаданных даёт значительные преимущества при работе с текстовыми документами, особенно в условиях роста объёмов информации и требований к её управлению \cite{gilliland2016metadata, dublincore}.

Улучшенный поиск становится возможным потому, что системы индексируют не только основной текст, но и метаданные. Это позволяет находить документы по автору, дате создания, ключевым словам, языку даже если в самом тексте эти слова не встречаются. Семантическая структура помогает поисковым алгоритмам лучше понимать релевантность документа запросу пользователя.

Автоматизация обработки документов упрощается, когда программы понимают их структуру. Можно автоматически генерировать оглавления, извлекать определённые типы контента (например, все цитаты или таблицы), преобразовывать документы в другие форматы с сохранением семантики. Это значительно экономит время и снижает вероятность ошибок при ручной обработке.

Совместимость между разными системами и форматами улучшается благодаря единым принципам организации. Документ, размеченный семантически, может быть преобразован из HTML в PDF или EPUB без потери смысловой структуры. Метаданные, хранящиеся в стандартизированном формате, могут быть прочитаны и использованы разными приложениями.

Долговременная сохранность информации обеспечивается тем, что семантическая разметка и метаданные хранятся в стандартизированном виде, который будет понятен будущим системам даже если конкретные программы устареют. Это особенно важно для архивов, библиотек и организаций, ответственных за сохранение цифрового наследия.

\ManualSubsection{Роль в современных информационных системах и будущее развитие}

В современных условиях семантическая организация и метаданные играют ключевую роль в системах управления контентом, электронных архивах, цифровых библиотеках и других информационных системах \cite{bernerslee2001semanticweb, gilliland2016metadata}.

Системы управления контентом (Content Management Systems, CMS) используют метаданные для категоризации, поиска и управления версиями документов. Семантическая разметка позволяет автоматически генерировать навигацию, связанные материалы, адаптивные представления для разных устройств. Это повышает эффективность работы с контентом и улучшает пользовательский опыт.

Электронные архивы и репозитории полагаются на метаданные для обеспечения долговременной сохранности и доступности цифровых объектов. Стандарты вроде Dublin Core определяют обязательный набор метаданных для архивного хранения. Семантическая организация контента позволяет создавать сложные связи между документами, обеспечивая целостность и контекст архивных коллекций.

Цифровые библиотеки используют семантическую организацию для создания интеллектуальных систем навигации и поиска. Метаданные позволяют связывать материалы, создавать тематические подборки, предлагать релевантный контент. Это превращает библиотеки из простых хранилищ в активные информационные среды.

Будущее развитие семантических технологий связано с развитием искусственного интеллекта и машинного обучения. Алгоритмы становятся способны самостоятельно определять структуру документа, идентифицировать типы контента, извлекать ключевые понятия и связи между ними. Семантические веб-технологии позволяют связывать документы в глобальную сеть знаний, где информация из разных источников может быть интегрирована и совместно использована.

Для эффективного использования семантической организации и метаданных рекомендуется использовать стандартные средства разметки, предоставляемые редакторами; заполнять метаданные максимально полно и точно; следовать принятым в конкретной области стандартам и соглашениям; периодически проверять и обновлять метаданные, особенно для долгоживущих документов.

Семантическая организация и метаданные превращают простые текстовые файлы в интеллектуальные информационные объекты, которые могут быть эффективно обработаны как людьми, так и программами. Понимание и применение этих принципов становится существенным навыком в цифровую эпоху, где объём информации растёт экспоненциально, а требования к её управлению и поиску постоянно повышаются.

\newpage
\refstepcounter{section}
\addcontentsline{toc}{section}{\protect\numberline{\thesection}Технологии автоматической обработки и анализа текстовой информации}
\begin{center}\textbf{\thesection\ ТЕХНОЛОГИИ АВТОМАТИЧЕСКОЙ ОБРАБОТКИ И АНАЛИЗА ТЕКСТОВОЙ ИНФОРМАЦИИ}\end{center}
\vspace*{-1.5em}

\ManualSubsection{Эволюция подходов: от правил к нейронным сетям}

Автоматическая обработка и анализ текстовой информации прошли значительную эволюцию, отражающую общее развитие компьютерных технологий и искусственного интеллекта. Каждый этап этой эволюции решал ограничения предыдущего, расширяя возможности и повышая точность обработки \cite{Lumenalta2025NLP, Coursera2025NLP}.

Начальный этап развития связан с системами на основе правил (rule-based systems). Эти системы представляли собой набор лингвистических правил и паттернов, составленных вручную разработчиками. Например, для распознавания имён собственных могли использоваться словари (gazetteers) или регулярные выражения, ищущие определённые последовательности символов. Преимущество этого подхода заключалось в полной прозрачности и контролируемости: каждое правило было понятно и объяснимо. Однако такие системы оказались крайне ограниченными в масштабируемости, требовали постоянного обновления правил для работы с новыми данными и языковыми явлениями, и не могли справиться с неоднозначностью, присущей естественному языку \cite{CEDAR2012IE}.

Следующим этапом стало внедрение статистических методов, основанных на машинном обучении. Модели, такие как скрытые марковские модели (Hidden Markov Models, HMM) и условные случайные поля (Conditional Random Fields, CRF), обучались на размеченных наборах данных (корпусах) и могли обобщать выявленные паттерны на новые тексты. Эти методы требовали подготовки обучающих данных, но затем могли автоматически выявлять закономерности без явного описания лингвистических правил. Точность обработки значительно возросла, системы стали более устойчивыми к вариациям в тексте \cite{GeeksforGeeks2024NER}.

Современный этап связан с глубоким обучением (deep learning) и нейронными сетями. Рекуррентные нейросети (RNN), особенно их модификации с долгой краткосрочной памятью (LSTM) и двунаправленные архитектуры (BiLSTM), способны учитывать контекстные зависимости в последовательностях текста. Настоящую революцию произвели трансформерные модели (Transformers), лежащие в основе таких систем, как BERT и GPT. Эти модели используют механизм внимания (attention mechanism), позволяющий одновременно обрабатывать все элементы текста и лучше понимать контекстное значение слов. Качественные улучшения очевидны: если системы на основе правил достигали точности 60–70\% на узких доменах, статистические методы обеспечивали 75–85\%, а современные трансформеры часто превышают 90\% точности на многих задачах \cite{DisplayR2025TextAnalysis, Lumenalta2025Tools}.

\ManualSubsection{Предварительная обработка текста: токенизация, стемминг, лемматизация}

Качество любого автоматического анализа текста в значительной степени зависит от этапа предварительной обработки. Этот этап включает несколько ключевых операций, которые подготавливают необработанный текст для последующего анализа \cite{Educative2023Tokenization, BecomingHuman2023NLP}.

Токенизация представляет собой процесс разбиения текста на отдельные элементы — токены. На первый взгляд эта задача кажется простой (разделить текст по пробелам), но на практике она сопряжена со значительными сложностями. Например, в предложении «Google's CEO is working on AI projects» необходимо правильно определить границы токенов: является ли «Google's» одним токеном или двумя («Google» и «'s»), как обрабатывать точки в аббревиатурах и многоточия. Грамотная токенизация учитывает языковые особенности и может использовать статистические методы или машинное обучение для определения оптимальных границ токенов \cite{GeeksforGeeks2024NER}.

Стемминг — это эвристический метод приведения словоформ к их основе (корню) путём отсечения аффиксов. Например, слова «walking», «walks», «walked» приводятся к общей форме «walk». Стемминг выполняется с помощью набора морфологических правил, специфичных для каждого языка. Этот метод отличается высокой скоростью работы и не требует знания словаря, но часто приводит к неточным результатам: «universe» может быть преобразовано в «univers», «argued» в «argu». Несмотря на эти недостатки, стемминг широко применяется в задачах информационного поиска, где скорость обработки важнее абсолютной точности \cite{Educative2023Tokenization}.

Лемматизация представляет собой более сложный и точный процесс приведения слов к их словарной форме (лемме). В отличие от стемминга, лемматизация учитывает часть речи слова и его морфологический анализ. Например, слово «better» лемматизируется в «good» (прилагательное в сравнительной степени), «running» — в «run» (глагол), «dances» — в «dance». Лемматизация требует использования специализированных словарей (например, WordNet) и POS-таггера (определителя части речи), поэтому она медленнее стемминга, но даёт качественно лучшие результаты. Этот метод незаменим для семантического анализа и задач, где критически важна точность определения смысла слов \cite{BecomingHuman2023NLP}.

Выбор конкретного метода предварительной обработки зависит от решаемой задачи. Для систем информационного поиска и поисковых движков часто достаточно стемминга благодаря его скорости. Для семантического анализа и глубокого понимания смысла текста необходима лемматизация. Современные нейросетевые модели иногда обходятся без явного стемминга или лемматизации, поскольку сами обучаются представлять синонимичные формы близко в семантическом пространстве \cite{Coursera2025NLP}.

\ManualSubsection{Распознавание именованных сущностей, синтаксический и семантический анализ}

Автоматический анализ текста включает несколько уровней, от выделения отдельных объектов до понимания структуры и смысла предложений \cite{GeeksforGeeks2024NER, SparkNLP2022DepParsing}.

Распознавание именованных сущностей (Named Entity Recognition, NER) решает задачу автоматического выявления и классификации реальных объектов, упоминаемых в тексте: имён людей, названий организаций, географических объектов, дат, сумм денег, научных терминов и других. Архитектура типичной NER-системы включает несколько этапов: токенизацию текста, POS-теггирование (присвоение частей речи), определение границ сущностей и их классификацию по заранее определённым категориям. Для реализации NER используются различные подходы: системы на основе знаний (словари и правила), статистическое машинное обучение (CRF) и глубокое обучение (BiLSTM, трансформеры). Современные модели достигают высокой точности, особенно в узких предметных областях \cite{Kairntech2025NER, arXiv2011NER}.

Синтаксический анализ (parsing) выстраивает грамматическую структуру предложения. Зависимостный парсинг (dependency parsing) определяет отношения между словами в предложении, назначая каждому слову голову (слово, от которого оно зависит синтаксически) и тип зависимости (подлежащее, дополнение, определение и др.). Алгоритмы парсинга делятся на переходные (пошагово строят дерево, применяя последовательность действий) и основанные на графах (вычисляют наиболее вероятное дерево для всего предложения). Синтаксический анализ важен для понимания грамматической структуры текста, что является основой для более сложных задач, таких как машинный перевод или генерация текста \cite{SparkNLP2022DepParsing}.

Семантический анализ идёт дальше синтаксиса, стремясь понять, что именно означает предложение. Этот уровень включает разрешение анафор (связывание местоимений с их антецедентами), извлечение отношений между сущностями (например, «работает в», «расположен в»), определение семантических ролей аргументов при предикате (кто что сделал кому) и понимание подразумеваемого смысла. Семантический анализ является одной из самых сложных задач обработки естественного языка, поскольку требует не только формального анализа структуры, но и привлечения фоновых знаний и понимания контекста. Современные подходы используют комбинацию машинного обучения, онтологий и больших языковых моделей \cite{ParallelDots2018NLP}.

Главное различие между этими уровнями анализа можно сформулировать так: NER отвечает на вопрос «Какие объекты упоминаются в тексте?», синтаксический анализ — «Как структурировано предложение?», а семантический анализ — «Что это предложение означает?».

\ManualSubsection{Анализ тональности и тематическое моделирование}

Два важных направления автоматического анализа текста — определение эмоциональной окраски и выявление скрытых тематических структур \cite{Techtarget2024Sentiment, JMLR2003LDA}.

Анализ тональности (sentiment analysis), также известный как анализ мнений (opinion mining), решает задачу определения эмоциональной окраски текста. Система должна определить, выражает ли автор текста положительное, отрицательное или нейтральное мнение об объекте обсуждения. Процесс анализа тональности обычно включает предварительную обработку текста (токенизацию, удаление стоп-слов), извлечение признаков (определение слов и фраз, коррелирующих с тональностью) и классификацию с использованием различных моделей: логистической регрессии, метода опорных векторов (SVM), наивного байесовского классификатора или нейронных сетей. Анализ может выполняться на разных уровнях: уровне документа (общая тональность всего текста), уровне предложения (тональность каждого предложения отдельно) и уровне аспектов (определение тональности относительно конкретных характеристик объекта, например, «еда вкусная, но сервис медленный»). Эта технология широко применяется для мониторинга брендов, анализа отзывов клиентов, изучения общественного мнения в социальных сетях \cite{GeeksforGeeks2024OpinionMining, AyaData2025OpinionMining}.

Тематическое моделирование (topic modeling) — это методология для автоматического обнаружения скрытых тематических структур в коллекциях документов. Наиболее известный метод — Latent Dirichlet Allocation (LDA). LDA представляет собой вероятностную модель, которая предполагает, что каждый документ является смесью нескольких тем, а каждая тема — это распределение вероятностей над словами. Математически модель использует байесовский вывод для «обратного инжиниринга» скрытых тематических структур на основе наблюдаемых слов в документах. Для получения апостериорного распределения применяются методы вроде семплирования Гиббса или вариационного байесовского вывода. Тематическое моделирование находит практическое применение в организации больших текстовых коллекций, анализе научных публикаций, изучении трендов в социальных медиа, категоризации новостей и автоматической кластеризации документов без явной разметки \cite{GeeksforGeeks2021LDA, Wikipedia2006LDA}.

\ManualSubsection{Автоматическое реферирование текстов}

Автоматическое реферирование (summarization) решает задачу создания краткого изложения исходного текста, сохраняющего ключевую информацию. Существуют два принципиально разных подхода к этой задаче \cite{GitHub2022Summarization, arXiv2023AbstractiveSummarization}.

Извлекающее реферирование (extractive summarization) работает по относительно простому принципу: из исходного текста отбираются те предложения, которые содержат наиболее важную информацию. Алгоритм обычно включает разбиение текста на предложения, оценку «важности» каждого предложения на основе различных признаков (частоты ключевых слов, позиции в документе, наличия значимых фраз) и выбор предложений с наивысшими оценками. Отобранные предложения затем объединяются, часто в исходном порядке. Преимущества этого подхода включают высокую скорость работы, относительную надёжность и отсутствие ошибок смысла, поскольку используются оригинальные выражения автора. Недостатки заключаются в возможной фрагментарности реферата, если выбранные предложения не связаны между собой естественно, и неспособности перефразировать или обобщать концепции.

Абстрактивное реферирование (abstractive summarization) представляет собой более сложную задачу: система генерирует новый текст, который передаёт ключевую информацию исходного документа, но может содержать фразы и конструкции, отсутствующие в оригинале. Этот подход ближе к тому, как человек пишет реферат: сначала понять смысл текста, а затем переизложить его своими словами. Реализация абстрактивного реферирования требует глубокого понимания текста (семантической компрехенции) и способности к генерации текста. Современные системы используют архитектуры типа seq2seq (sequence-to-sequence) на основе трансформеров, такие как BART или T5, обученные на больших корпусах пар «текст-реферат». Преимущества абстрактивного подхода включают более естественное звучание реферата, способность обобщать и перефразировать концепции. Недостатки — высокие требования к вычислительным ресурсам, риск «галлюцинаций» (генерации информации, которой нет в оригинале) и зависимость от качества обучающих данных.

На практике часто применяются гибридные подходы, которые комбинируют извлекающее и абстрактивное реферирование. Например, сначала извлекаются ключевые предложения, которые затем передаются в абстрактивную модель для генерации более связного и естественного реферата. Автоматическое реферирование находит применение в создании дайджестов новостей, аннотировании научных статей, подготовке обзоров документов и многих других областях, где требуется быстро получить суть больших объёмов текстовой информации \cite{MicrosoftLearn2025Summarization}.

Развитие технологий автоматической обработки и анализа текста продолжает ускоряться благодаря прогрессу в области искусственного интеллекта, увеличению вычислительных мощностей и доступности больших текстовых данных. Эти технологии уже сегодня трансформируют работу с информацией в бизнесе, науке, образовании и повседневной жизни, открывая новые возможности для извлечения знаний из текстовых данных.

\newpage

\refstepcounter{section}
\addcontentsline{toc}{section}{\protect\numberline{\thesection}Основы долговременного хранения и архивации текстовых ресурсов}
\begin{center}\textbf{\thesection\ ОСНОВЫ ДОЛГОВРЕМЕННОГО ХРАНЕНИЯ И АРХИВАЦИИ ТЕКСТОВЫХ РЕСУРСОВ}\end{center}
\vspace*{-1.5em}

\ManualSubsection{Стратегии долговременного сохранения: миграция, обновление, эмуляция}

Цифровые технологии развиваются стремительно: аппаратное обеспечение устаревает, программное обеспечение становится недоступным, форматы файлов выходят из употребления. Для обеспечения доступности текстовых ресурсов на протяжении десятилетий и даже столетий необходимы специальные стратегии долговременного сохранения \cite{Lucidea2024LongTermAccess, ETHLibrary2025DigitalPreservation}.

Обновление носителя (refreshment) представляет собой самую простую стратегию. Данные физически не изменяются и не конвертируются, но периодически (каждые 5–10 лет) копируются со старых, деградирующихся носителей информации на новые. Например, данные с магнитных лент LTO третьего поколения переносятся на ленты пятого поколения до того, как старые носители выйдут из строя. Преимущества этого подхода включают простоту реализации, отсутствие необходимости изменять сами данные и относительно низкую стоимость. Однако обновление носителя не решает проблему устаревания форматов файлов, требует регулярного мониторинга и в конечном итоге может потребовать перехода на носители, которые больше не поддерживаются производителями \cite{UniversityOfKentucky1994MassStorage}.

Миграция формата (format migration) — это перевод данных из одного формата в другой, более современный или лучше подходящий для долговременного хранения. Например, документ в старом формате .doc (Microsoft Word 97-2003) может быть конвертирован в стандартизированный формат PDF/A. Этот подход непосредственно решает проблему форматного устаревания. Существует несколько типов миграции: миграция версии (обновление документа на новую версию того же формата), миграция программного обеспечения (переход на другое приложение с похожей поддержкой) и полная миграция формата. Процесс включает определение целевого формата, использование инструментов конвертации, проверку целостности и точности конвертированных данных и сохранение метаданных о процессе миграции. Хотя миграция эффективно решает проблему совместимости, она сопряжена с риском потери данных или функциональности, требует значительных трудозатрат и экспертизы, а также может потребовать многократного повторения (каскадная миграция) \cite{SmithsonianInstitution2017PreservationFormats, Cornell2025ArchivingFileFormats}.

Эмуляция (emulation) предлагает принципиально иной подход. Вместо изменения самих данных создаётся виртуальное окружение, которое имитирует оригинальное аппаратное обеспечение, операционную систему и программное обеспечение. Например, для просмотра документа, созданного в WordPerfect 5.1 на компьютере Commodore 64, можно использовать эмулятор этой системы, работающий на современном компьютере. Эмуляция сохраняет исходный «вид и ощущение» документа, оригинальные данные остаются неизменными, и этот подход может применяться к огромному диапазону цифровых объектов одновременно. Однако эмуляция требует значительного технического мастерства и ресурсов, необходимо создавать эмулятор для каждой комбинации ОС и ПО, возникают вопросы интеллектуальной собственности, и неясно, будут ли будущие системы поддерживать современные эмуляторы. В долгосрочной перспективе эмуляция может оказаться очень дорогостоящей стратегией \cite{LIS654DigitalPreservation2015MigrationVsEmulation, DocByte2024PreservationStrategies}.

На практике организации часто применяют гибридный подход, комбинируя различные стратегии. Например, критически важные документы могут сохраняться с помощью эмуляции для сохранения оригинального внешнего вида, одновременно создаётся мигрированная версия для обеспечения доступности, и используется обновление носителей для физической сохранности данных \cite{NationalArchivesUK2025PreservationTools}.

\ManualSubsection{Форматы долговременного хранения: PDF/A, XML и открытые стандарты}

Выбор формата хранения является одним из самых критических решений при создании цифрового архива. Формат должен быть стабильным, хорошо задокументированным и желательно открытым, чтобы гарантировать доступность данных в будущем \cite{PARADISEC2024StandardFormats, EmersonCollege2025PreservationFormats}.

PDF/A (Portable Document Format for Archiving) — это семейство стандартов, определённых ISO 19005 специально для долговременного архивирования электронных документов. История развития PDF/A включает несколько версий: PDF/A-1 (основан на PDF 1.4, 2005), PDF/A-2 (на основе PDF 1.7, 2011, позволяет встраивать совместимые вложения), PDF/A-3 (2012, разрешает встраивание файлов в любом формате) и развивающийся PDF/A-4 (на основе PDF 2.0). Ключевые характеристики PDF/A включают самодостаточность (все необходимые шрифты и ресурсы встраиваются в файл), детерминированность внешнего вида (запрещены интерактивные и динамические элементы), открытость и стандартизацию (спецификация является открытым международным стандартом). PDF/A определяет несколько уровней соответствия: Level A (Accessible, требует полной доступности), Level B (Basic, фокусируется на воспроизведении визуального вида) и Level U (Unicode, обеспечивает извлечение текста) \cite{LOC2024PDFA, PDFAssociation2025ISO19005, PDFAssociation2025ISO19005Part1}.

XML (eXtensible Markup Language) широко рекомендуется для долговременного архивирования структурированных текстовых данных. Преимущества XML для архивации включают его текстовую природу (не бинарный формат), что обеспечивает долгосрочную читаемость; иерархическую структуру, хорошо подходящую для представления сложных документов; статус открытого стандарта W3C; лёгкость парсинга и обработки даже примитивными инструментами; поддержку Unicode для многоязычных текстов; возможность дополнения XML Schema для описания валидной структуры. XML особенно эффективен для хранения документов с чёткой структурой, таких как законодательные акты, научные статьи, каталоги, метаданные \cite{Sphereon2020DigitalPreservationPDFA, DigitalPreservation2008OfficeOpenXMLStandards}.

Открытые форматы в целом предпочтительнее проприетарных для долговременного хранения. Проприетарные форматы (например, .docx, .xlsx) представляют значительный риск: производитель может прекратить поддержку формата, спецификация может измениться без предупреждения, требуется лицензионное программное обеспечение для доступа, а сложность формата затрудняет миграцию. Открытые, стандартизированные форматы, такие как PDF/A, XML, простой текст в UTF-8, обеспечивают гораздо более высокий уровень уверенности в долговременной доступности данных. Для различных типов текстовых ресурсов существуют рекомендуемые форматы: для текстовых документов — PDF/A-1, plain text (UTF-8); для структурированных данных — XML; для веб-контента — WARC, ARC; для электронной корреспонденции — EML \cite{Soutron2025TrustedDigitalRepository, Wikipedia2004ArchiveFormats}.

\ManualSubsection{Проверка целостности данных: контрольные суммы и фиксити}

Одна из наиболее серьёзных угроз для архивированных данных — неявная коррупция, когда файл повреждается, но это остаётся незамеченным до момента попытки его использования. Такое повреждение может быть вызвано ошибками аппаратного обеспечения, сбоями в системах хранения, битовым сгниванием (bit rot) или воздействием вредоносного программного обеспечения \cite{OrbisFileFixity2014, LibraryOfCongress2012FileFixityDataIntegrity}.

Целостность данных (fixity) — это уверенность в том, что цифровой объект не был изменён с момента его поступления в архив, то есть остался «фиксированным». Обеспечение и проверка целостности достигаются с помощью криптографических хеш-функций, которые создают уникальные цифровые отпечатки содержимого файлов \cite{OrbisFileFixityDetailed2008}.

Контрольная сумма (checksum) или хеш-сумма представляет собой результат математической функции, преобразующей всё содержимое файла в строку фиксированной длины. Ключевое свойство криптографических хеш-функций заключается в том, что даже минимальное изменение одного бита в исходном файле приводит к совершенно другому значению контрольной суммы. Это позволяет обнаруживать любые, даже самые незначительные, изменения в данных \cite{LibraryOfCongress2012FileFixityDataIntegrity}.

Наиболее используемыми алгоритмами являются MD5 (128-битный хеш, сейчас не рекомендуется для критических приложений из-за возможности коллизий), SHA-1 (160-битный, также постепенно выводится из использования), SHA-256 (часть семейства SHA-2, 256-битный, рекомендуется для новых систем) и BLAKE2 (современный быстрый алгоритм). Для долговременных архивов рекомендуется использование алгоритмов семейства SHA-2 или SHA-3 ввиду их криптографической стойкости \cite{OrbisFileFixityDetailed2008}.

Процесс обеспечения целостности данных в архиве включает несколько этапов. При поступлении файла (ингезте) вычисляется его контрольная сумма, которая сохраняется вместе с метаданными объекта. Через определённые периодические интервалы (например, ежегодно) система перевычисляет контрольные суммы всех хранимых файлов и сравнивает их с исходными значениями. Если обнаружено несовпадение, это сигнализирует о повреждении файла. Архивы обычно хранят несколько копий данных (репликация) или используют коды исправления ошибок, что позволяет восстановить повреждённую копию из неповреждённой. Уровни реализации могут варьироваться от базового (генерация контрольных сумм при ингезте) до продвинутого (регулярная автоматическая проверка, система самовосстановления) и высокого (интеграция с файловыми системами, поддерживающими целостность данных, такими как ZFS или CEPH) \cite{ElectronicArtsIntermix1996PreservationStrategies}.

\ManualSubsection{Метаданные и стандарт Dublin Core в контексте архивации}

Метаданные — данные о данных — играют критически важную роль в долговременном сохранении цифровых объектов. Для цифрового архива метаданные обеспечивают идентификацию и описание ресурса, контроль доступа и прав, возможность восстановления информации в будущем и подтверждение подлинности объекта \cite{DublinCorePapers2005MetadataApproach, ZB_MED2016MetadataPreservation}.

Dublin Core — это простой, но мощный стандарт описания метаданных для цифровых ресурсов, разработанный на конференции в Дублине (Огайо) в 1995 году и впоследствии стандартизированный как ISO 15836. Стандарт определяет 15 основных элементов, каждый из которых описывает определённый аспект цифрового объекта \cite{DublinCore, OrangeLogic2024DublinCore}.

К этим элементам относятся: Title (название), Creator (создатель), Subject (тема, ключевые слова), Description (описание), Publisher (издатель), Contributor (соавтор), Date (дата), Type (тип ресурса), Format (формат файла), Identifier (уникальный идентификатор), Source (источник), Language (язык), Relation (связь с другими ресурсами), Coverage (пространственный или временной охват), Rights (права и лицензирование). Dublin Core также поддерживает квалифицированные расширения, позволяющие уточнять элементы через свойства (например, date.created, date.modified) и схемы кодирования \cite{YouTube2023DublinCoreOmeka}.

В контексте архивного хранения Dublin Core применяется для описания текстовых ресурсов, обеспечивая их надёжную идентификацию и поиск как сейчас, так и в будущем. Например, для архивированного письма могут быть указаны автор, дата написания, адресат, язык, тематика, уникальный идентификатор и условия доступа. Без таких метаданных будущие исследователи могут оказаться неспособными понять происхождение, контекст и назначение сохранённого документа, что сделает его бесполезным \cite{DublinCorePapers2005MetadataApproach}.

\ManualSubsection{Баланс между сохраняемостью и доступностью}

В долговременном сохранении цифровых объектов существует тонкий, но критически важный баланс между двумя концепциями: сохраняемостью (preservation) и доступностью (accessibility). Эти цели часто конкурируют друг с другом, и их оптимальное соотношение является ключевой задачей при проектировании архивных систем \cite{Lucidea2024LongTermAccess, ETHLibrary2025DigitalPreservation}.

Сохраняемость фокусируется на долговременном поддержании целостности, аутентичности и неизменности цифровых объектов. Это включает гарантию того, что данные не повреждены, сохранение истории и контекста объекта, возможность доказать его подлинность и защиту от несанкционированных изменений. Меры сохраняемости включают хранение в неизменяемых форматах, создание множественных копий на разных физических локациях, регулярную проверку целостности и логирование всех операций с объектом (audit trail) \cite{Rutgers2005ArchitectureTrustedDigitalRepositories}.

Доступность сосредоточена на способности текущих и будущих пользователей фактически прочитать и использовать архивированный контент. Это означает, что документ должен быть открываем и читаем, должны быть доступны необходимые инструменты и программное обеспечение, контент должен быть понятен и пригоден к использованию, а права доступа должны быть ясны. Меры доступности включают конвертирование в популярные, легко читаемые форматы, создание нескольких версий одного документа, предоставление онлайн-интерфейсов для доступа, обеспечение полнотекстового поиска и трансформацию в веб-дружественные форматы \cite{Soutron2025TrustedDigitalRepository}.

Конфликты между сохраняемостью и доступностью проявляются в нескольких аспектах. Для максимальной сохраняемости лучше хранить документы в простых, неизменяемых форматах (XML, PDF/A), но пользователи часто ожидают получить документ в более привычных или функциональных форматах. Для сохраняемости нежелательна любая модификация оригинального файла, но для доступности иногда необходимо обновлять или конвертировать формат. Полные метаданные необходимы для сохраняемости, но могут усложнять простой доступ. Создание множественных копий улучшает сохраняемость, но увеличивает риск расхождения версий \cite{LIS654DigitalPreservation2015MigrationVsEmulation}.

Оптимальным решением этого противоречия является принцип «консервации плюс доступность», реализуемый через модель OAIS (Open Archival Information System), стандартизированную как ISO 14721:2012. Эта модель предполагает создание отдельных, но связанных версий документа: Архивного информационного пакета (Archival Information Package, AIP), который хранится в самом консервативном формате с полными метаданными; Пакета распространения информации (Dissemination Information Package, DIP) — конвертированного в удобные форматы для пользователей; и Промежуточного информационного пакета (Submission Information Package, SIP) — формата, в котором информация поступает в архив. Такой подход обеспечивает баланс между долговременной сохраняемостью и текущей доступностью цифровых ресурсов \cite{Rutgers2005ArchitectureTrustedDigitalRepositories}.

Долговременное хранение и архивация текстовых ресурсов представляют собой комплексную междисциплинарную задачу, требующую сочетания технических, организационных и нормативных подходов. Понимание основных принципов и стратегий архивации необходимо для обеспечения сохранности цифрового культурного, научного и административного наследия для будущих поколений.


\newpage
\section*{}
\vspace*{-2.5em}
\begin{center}\textbf{ЗАКЛЮЧЕНИЕ}\end{center}
\phantomsection
\addcontentsline{toc}{section}{Заключение}

Технологии создания и обработки текстовых документов образуют целостный и логически связанный комплекс методов и инструментов, охватывающий весь жизненный цикл работы с текстовой информацией — от её фундаментального представления в цифровой форме до сложного анализа и долговременного хранения. Исследование показало, что эффективность работы с текстовыми данными определяется не только владением конкретными программными средствами, но и глубоким пониманием системных принципов, лежащих в основе каждого этапа этого процесса.

Фундаментом всей цифровой текстовой информации является система кодирования. Эволюция от ограниченного ASCII через хаос несовместимых 8-битных кодировок к универсальному стандарту Unicode и эффективным алгоритмам кодирования UTF, особенно UTF-8, обеспечила возможность корректного представления текстов на всех языках мира. Понимание принципов кодирования критически важно для предотвращения распространённых проблем с отображением текста и обеспечения интернационализации приложений.

Выбор формата хранения текстовых документов должен быть осознанным, учитывающим баланс между функциональностью, совместимостью и долговременной сохраняемостью. Открытые стандартизированные форматы, такие как ODT, PDF/A, XML, обеспечивают независимость от конкретных вендоров и гарантируют доступность данных в будущем, в то время как проприетарные форматы могут предлагать расширенные возможности в рамках замкнутых экосистем. Осознанный выбор формата является важным элементом цифровой грамотности и ответственного управления информацией.

Языки разметки играют ключевую роль в структурировании документов, отделяя содержание от оформления и обеспечивая машиночитаемость текста. От HTML и Markdown до LaTeX и XML — эти инструменты позволяют создавать документы, которые легко обрабатывать автоматически, преобразовывать в различные форматы и сохранять на долгие годы. Семантическая организация и метаданные превращают простые текстовые файлы в интеллектуальные информационные объекты, обогащая их смысловой структурой и контекстуальной информацией, что значительно повышает возможности их автоматической обработки, поиска и интеграции.

Современные технологии автоматической обработки и анализа текста, от предварительной обработки до глубокого семантического анализа, открывают новые возможности для извлечения знаний из текстовых данных. Распознавание именованных сущностей, анализ тональности, тематическое моделирование, автоматическое реферирование — эти и другие методы позволяют преобразовывать неструктурированный текст в структурированную информацию, пригодную для анализа и принятия решений. Развитие нейросетевых моделей, особенно трансформеров, продолжает расширять границы возможного в понимании естественного языка.

Долговременное хранение и архивация текстовых ресурсов требуют системного подхода, сочетающего технические стратегии (миграцию, эмуляцию), использование специализированных форматов (PDF/A, XML), обеспечение целостности данных (контрольные суммы) и ведение подробных метаданных (Dublin Core). Баланс между сохраняемостью и доступностью, реализуемый через модели вроде OAIS, обеспечивает как защиту цифрового наследия для будущих поколений, так и его практическую полезность в настоящем.

Таким образом, технологии работы с текстовыми документами представляют собой взаимосвязанную экосистему, охватывающую все аспекты жизненного цикла текстовой информации. Комплексное владение этой экосистемой — от понимания принципов кодирования и выбора форматов до применения продвинутых методов анализа и стратегий архивации — формирует ключевую компетенцию специалиста в современном цифровом обществе. Это позволяет не только эффективно решать текущие задачи, но и создавать масштабируемые, воспроизводимые и надёжные процессы работы с текстовой информацией, что является неотъемлемым условием успеха в условиях продолжающейся цифровой трансформации всех сфер человеческой деятельности.

\newpage
\begin{center}\textbf{СПИСОК СОКРАЩЕНИЙ}\end{center}
\phantomsection
\addcontentsline{toc}{section}{Список сокращений}

\vspace{1em}

{\small
\noindent
\begin{tabular}{|p{0.15\textwidth}|p{0.80\textwidth}|}
\hline
\textbf{Сокращение} & \multicolumn{1}{c|}{\textbf{Расшифровка}} \\
\noalign{\hrule height 1.2pt}
\hline
ASCII & American Standard Code for Information Interchange (Американский стандартный код для обмена информацией) \\
\hline
UTF & Unicode Transformation Format (Формат преобразования Юникода) \\
\hline
HTML & HyperText Markup Language (Язык гипертекстовой разметки) \\
\hline
XML & eXtensible Markup Language (Расширяемый язык разметки) \\
\hline
JSON & JavaScript Object Notation (Обозначение объектов JavaScript) \\
\hline
PDF & Portable Document Format (Портативный формат документа) \\
\hline
PDF/A & Portable Document Format for Archiving (Портативный формат документа для архивирования) \\
\hline
ODF & Open Document Format (Открытый формат документов) \\
\hline
OOXML & Office Open XML (Открытый XML для офиса) \\
\hline
CSV & Comma-Separated Values (Значения, разделённые запятыми) \\
\hline
NLP & Natural Language Processing (Обработка естественного языка) \\
\hline
NER & Named Entity Recognition (Распознавание именованных сущностей) \\
\hline
LDA & Latent Dirichlet Allocation (Латентное размещение Дирихле) \\
\hline
BERT & Bidirectional Encoder Representations from Transformers (Двунаправленные представления энкодера от трансформеров) \\
\hline
GPT & Generative Pre-trained Transformer (Генеративный предварительно обученный трансформер) \\
\hline
RNN & Recurrent Neural Network (Рекуррентная нейронная сеть) \\
\hline
LSTM & Long Short-Term Memory (Долгая краткосрочная память) \\
\hline
CRF & Conditional Random Fields (Условные случайные поля) \\
\hline
SVM & Support Vector Machine (Метод опорных векторов) \\
\hline
CMS & Content Management System (Система управления контентом) \\
\hline
OAIS & Open Archival Information System (Открытая архивная информационная система) \\
\hline
BOM & Byte Order Mark (Метка порядка байтов) \\
\hline
POS & Part of Speech (Часть речи) \\
\hline
W3C & World Wide Web Consortium (Консорциум Всемирной паутины) \\
\hline
ISO & International Organization for Standardization (Международная организация по стандартизации) \\
\hline
\end{tabular}
}

\clearpage
\nocite{*}
\printbibliography[heading=bibliography]

\end{document}
